\subsection{SEMI-DIRECT PRODUCTS OF LIE GROUPS}

Let I be a finite set and \((\mathbf{L}_{i})_{i\in I}\) a family of Lie groups. The group and manifold structures on \(L=\prod_{i\in I}L_{i}\) are compatible and L thus has a Lie group structure. L is called the product Lie group of the family of Lie groups \((\mathbf{L}_{i})_{i\in I}\) .

Let L and M be Lie groups and \(\sigma\) a homomorphism of L into the automorphism group of the group M. Let S be the external semi-direct product of L by M relative to \(\sigma\) (Algebra, Chapter I, §6, no. 1, Definition 2).

PROPOSITION 7. If the mapping \((m, l) \mapsto \sigma(l)m\) of \(\mathbf{M} \times \mathbf{L}\) into \(\mathbf{M}\) is analytic, the group \(\mathbf{S}\) , with the product manifold structure of \(\mathbf{M}\) and \(\mathbf{L}\) , is a Lie group.

For \(l, l'\) in L and \(m, m'\) in M,

\[
\begin{array}{r l} (m, l) (m ^ {\prime}, l ^ {\prime}) ^ {- 1} & = m l l ^ {- 1} m ^ {\prime - 1} = m (\sigma (l l ^ {\prime - 1}) m ^ {\prime - 1}) l l ^ {\prime - 1} \\ & = (m (\sigma (l l ^ {\prime - 1}) m ^ {\prime - 1}), l l ^ {\prime - 1}) \end{array}
\]

whence the proposition.

If the conditions of Proposition 7 hold, the Lie group S is called the (external) semi-direct product Lie group of L by M relative to \(\sigma\) .

Clearly the canonical injection of L (resp. M) into S is an isomorphism of L (resp. M) onto a Lie subgroup of S which we identify with L (resp. M). The canonical mapping of S onto L is a Lie group morphism.

Conversely, let G be a Lie group and L, M two Lie subgroups such that the group G is (algebraically) the semi-direct product of L by M(Algebra, Chapter I, § 6, no. 1). We write \(\sigma(l)m = lml^{-1}\) for \(l \in L\) and \(m \in M\) . Then \(\sigma\) satisfies the conditions of Proposition 7. We can therefore form the semi-direct product Lie group S of L by M relative to \(\sigma\) . The mapping \(j: (m, l) \mapsto ml\) of S onto G is a group isomorphism and is analytic. If \(j\) is a Lie group isomorphism, the Lie group G is called the (internal) semi-direct product of L by M and S and G are identified. For all \(g \in G\) , we write \(g = p(g)q(g)\) , where \(p(g) \in M\) and \(q(g) \in L\) . For the Lie group G to be the semi-direct product of L by M, it is necessary and sufficient that one of the mappings \(p: G \to M\) and \(q: G \to L\) be analytic, in which case both are analytic; or alternatively, it is necessary and sufficient that \(\mathrm{T}_{e}(\mathrm{G})\) be the topological direct sum of \(\mathrm{T}_{e}(\mathrm{M})\) and \(\mathrm{T}_{e}(\mathrm{L})\) (for, if this condition holds, \(j\) is étale at \(e_{\mathrm{s}}\) ).

Example. Let E be a normable space, \(G = \mathbf{GL}(E)\) , T the translation group of E and A the permutation group of E generated by G and T. The group A is algebraically the semi-direct product of G by T. (If E is finite-dimensional, A is the affine group of E, cf.~Algebra, Chapter II, § 9, no. 4). Let \(\sigma\) be the identity linear representation of G on E and S the external semi-direct product of G by E relative to \(\sigma\) . For all \(x \in E\) , let \(t_{x}\) be the translation of E defined by x. The mapping \((x, u) \mapsto t_{x} \circ u\) is isomorphism \(\Phi\) of the group S onto the group A. The mapping \((x, u) \mapsto \sigma(u)x = u(x)\) of \(E \times \mathcal{L}(E)\) into E is continuous and bilinear and hence analytic; its restriction to \(E \times G\) is therefore analytic. Thus the group S, with the product manifold structure of E and G, is a Lie group. We transport this structure to A by means of \(\Phi\) . Then A becomes a Lie group, the internal semi-direct product of G by T as a Lie group.

PROPOSITION 8. Let G and H be Lie groups, \(p: \mathrm{G} \to \mathrm{H}\) and \(s: \mathrm{H} \to \mathrm{G}\) Lie group morphisms such that \(p \circ s = \mathrm{id}_{\mathrm{H}}\) and \(N = \operatorname{Ker} p\) . Then \(N\) is a Lie subgroup of \(G\) , \(s\) is an isomorphism of \(H\) onto a Lie subgroup of \(G\) and the Lie group \(G\) is the internal semi-direct product of \(s(H)\) by \(N\) .

\(T_{e}(p) \circ T_{e}(s) = \mathrm{id}_{T_{e}(\mathrm{H})}\) and hence \(p\) (resp. \(s\) ) is a submersion (resp. an immersion). By Differentiable and Analytic Manifolds, R, 5.10.5, N is a Lie subgroup of G. On the other hand, \(s\) is a homeomorphism of H onto \(s(H)\) and hence \(s\) is an isomorphism of H onto a Lie subgroup of G (Differentiable and Analytic Manifolds, R, 5.8.3). Finally, for all \(g \in G\) , \(g = (s \circ p)(g)\) . \(n\) for some \(n \in N\) ; as \(s \circ p\) is analytic, the Lie group G is the semi-direct product of \(s(H)\) by N.

